\section{Preliminaries}\label{sec:prep}

We start with f\/ixing the notations.

\subsection{Notations}\label{sec:notations}

Manifolds are always assumed to be smooth.
If $M$, $N$ are manifolds and $f\colon M\rightarrow N$ is a~smooth map, then $\mathrm{d} f\colon TM\rightarrow TN$ denotes
the dif\/ferential map between their tangent manifolds.
The map $f$ is said to be an immersion if\/f for each $x\in M$ the restriction $\mathrm{d}_xf:=\mathrm{d} f|_{T_xM}\colon
T_xM\rightarrow T_{f(x)}N$ is injective.

Let $V$ be a~f\/inite dimensional vector space.
A $V$-valued 1-form $\omega$ on the manifold $N$ is a~smooth map $\omega\colon TN\rightarrow V$ whose restriction
$\omega_y:=\omega|_{T_yN}$ is linear for all $y\in N$.
The pullback of $\omega$ by $f$ is the $V$-valued 1-form $f^{*}\omega\colon TM\rightarrow V$, $\vec{v}_x\rightarrow
\omega_{f(x)}(\mathrm{d}_xf(\vec{v}_x))$.

Let~$G$ be a~Lie group with Lie algebra $\mathfrak{g}$.
For $g\in G$, we def\/ine the corresponding conjugation map by $\Co{g}\colon G\rightarrow G$, $h\mapsto g h g^{-1}$.
Its dif\/ferential $\mathrm{d}_e\Co{g}\colon \mathfrak{g}\rightarrow \mathfrak{g}$ at the unit element $e\in G$ is denoted
by $\Add{g}$ in the following.

Let $\Psi$ be a~(left) action of the Lie group~$G$ on the manifold~$M$.
For $g\in G$ and $x\in M$, we define $\Psi_g\colon M\rightarrow M$, $\Psi_g\colon y\mapsto \Psi(g,y)$ and $\Psi_x\colon G\rightarrow M$, $h\mapsto\Psi(h,x)$, respectively. If it is clear which action is meant, we will often write $L_g$ instead of $\Psi_g$ as well as $g\cdot y$ or $g y$ instead of $\Psi_g(y)$. For $\vec{g}\in \mathfrak{g}$ and $x\in M$, the map
\begin{gather*}
\wt{g}(x):=\dttB{t}{0} \Psi_x(\exp(t\vec{g}\hspace{1pt}))
\end{gather*}
is called the
\emph{fundamental vector field of~$\vec{g}$}.
The Lie subgroup $G_x:=\left\{g\in G\: \big| \: g\cdot x=x\right\}$ is called the \emph{stabilizer} of $x\in M$
(w.r.t.~$\Psi$), and its Lie algebra $\mathfrak{g}_x$ equals $\ker[\mathrm{d}_x\Psi]$, see e.g.~\cite{DuisKolk}.
The \emph{orbit} of $x$ under~$G$ is the set $Gx:=\operatorname{\mathrm{im}}[\Psi_x]$. $\Psi$ is said to be
\emph{transitive} if\/f $Gx=M$ holds for one (and then each) $x\in M$.
Analogous conventions we also use for right actions.

\subsection{Invariant connections}\label{subsec:InvConn}

Let $\pi\colon P\rightarrow M$ be a~smooth map between manifolds $P$ and~$M$, and denote by
$F_x:=\pi^{-1}(x)\subseteq P$ the f\/ibre over $x\in M$ in $P$.
Moreover, let $S$ be a a~Lie group that acts via $R\colon P\times S\rightarrow P$ from the right on $P$.
If there is an open covering $\{U_\alpha\}_{\alpha\in I}$ of~$M$ and a~family $\{\phi_\alpha\}_{\alpha\in I}$ of
dif\/feomorphisms $\phi_\alpha\colon \pi^{-1}(U_\alpha)\rightarrow U_\alpha\times S$ with
\begin{gather}
\label{eq:bundlemaps}
\phi_\alpha(p\cdot s)=\big(\pi(p),[\mathrm{pr}_2\circ\phi_\alpha](p)\cdot s\big)
\qquad\forall\, p\in \pi^{-1}(U_\alpha),\qquad\forall\, s\in S,
\end{gather}
then $(P,\pi,M,S)$ is called \emph{principal fibre bundle} with total space $P$,
projection map $\pi$, base manifold~$M$ and structure group $S$.
Here, $\mathrm{pr}_2$
denotes the projection onto the second factor.
It follows from~\eqref{eq:bundlemaps} that $\pi$ is surjective, and that:
\begin{itemize}\itemsep=0pt
\item
$R_s(F_x)\subseteq F_x$ for all $x\in M$ and all $s\in S$,
\item
for each $x\in M$ the map $R_x\colon F_x \times S\rightarrow F_x$, $(p,s)\mapsto p\cdot s$ is transitive and free.
\end{itemize}
The subspace $Tv_pP:=\ker[d_p\pi]\subseteq T_pP$ is called \emph{vertical tangent
space} at $p\in P$ and
\begin{gather*}
\wt{s}(p):=\dttB{t}{0}\: p\cdot \exp(t\vec{s}\hspace{1pt})\in Tv_pP
\qquad\forall\, p\in P,
\end{gather*}
denotes the fundamental vector f\/ield of $\vec{s}$ w.r.t.~the right action of $S$ on
$P$.
The map $\mathfrak{s}\ni\vec{s}\rightarrow \wt{s}(p)\in Tv_pP$ is a~vector space isomorphism for all $p\in P$.

Complementary to that, a~\emph{connection} $\omega$ is an $\mathfrak{s}$-valued 1-form on $P$ with \begin{itemize}\itemsep=0pt
\item
$R_s^{*}\omega= \Add{s^{-1}}\circ \omega\qquad \forall\,s\in S$,
\item
$\omega_p(\wt{s}(p))=\vec{s}\hspace{44pt}\forall\,\vec{s}\in \mathfrak{s}$.
\end{itemize}
The subspace $Th_pP:=\ker[\omega_p]\subseteq T_pP$ is called the \emph{horizontal
tangent space} at $p$ (w.r.t.~$\omega$).
We have $\mathrm{d} R_s(Th_pP)=Th_{p\cdot s}P$ for all $s\in S$, and one can show that $T_pP= Tv_pP\oplus Th_pP$ holds for all
$p\in P$.

A dif\/feomorphism $\kappa\colon P\rightarrow P$ is said to be an \emph{automorphism} if\/f $\kappa(p\cdot s)=\kappa(p)\cdot
s$ holds for all $p\in P$ and all $s\in S$.
It is straightforward to see that an $\mathfrak{s}$-valued 1-form $\omega$ on $P$ is a~connection if\/f this is true for
the pullback $\kappa^{*}\omega$.
A \emph{Lie group of automorphisms $(G,\Phi)$ of $P$} is a~Lie group~$G$ together with a~left action $\Phi$ of~$G$ on
$P$ such that the map $\Phi_g$ is an automorphism for each $g\in G$.
This is equivalent to say that $\Phi(g,p\cdot s)=\Phi(g,p)\cdot s$ holds for all $p\in P$, $g\in G$ and all $s\in S$.
In this situation, we will often write $gps$ instead of $(g\cdot p)\cdot s=g\cdot(p\cdot s)$.
Each such a~left action $\Phi$ gives rise to two further actions: \begin{itemize}\itemsep=0pt
\item
The induced action $\varphi$ is def\/ined by
\begin{gather}
\label{eq:INDA}
\begin{split}
&\varphi\colon\quad G\times M  \rightarrow  M,
\\
& \phantom{\varphi\colon\quad{}}
(g,m) \mapsto  (\pi\circ\Phi)(g, p_m),
\end{split}
\end{gather}
where $p_m\in \pi^{-1}(m)$ is arbitrary.
$\Phi$ is called \emph{fibre transitive} if\/f $\varphi$ is transitive.
\item
We equip $Q=G\times S$ with the canonical Lie group structure and def\/ine~\cite{Wang}
\begin{gather}
\label{eq:THETA}
\begin{split}
&\Theta\colon\quad Q \times P \rightarrow  P,
\\
&\phantom{\Theta\colon\quad{}}
((g,s),p) \mapsto  \Phi\left(g, p\cdot s^{-1}\right).
\end{split}
\end{gather}
\end{itemize}
A connection $\omega$ is said to be $\Phi$-invariant if\/f $\Phi_g^{*}\omega=\omega$ holds for all
$g\in G$.
This is equivalent to require that for each $p\in P$ and $g\in G$ the dif\/ferential $\mathrm{d}_pL_g$ induces an
isomorphism between the horizontal tangent spaces $Th_pP$ and $Th_{gp}P$.\footnote{In literature sometimes the latter
condition is used to def\/ine $\Phi$-invariance of connections.}

We conclude this subsection with the following straightforward facts, see also~\cite{Wang}:
\begin{itemize}\itemsep=0pt
\item
Consider the representation $\rho\colon Q \rightarrow \mathsf{Aut}(\mathfrak{s})$, $(g,s)\mapsto \Add{s}$.
Then, it is straightforward to see that each $\Phi$-invariant connection $\omega$ is of type $\rho$, i.e., $\omega$ is an
$\mathfrak{s}$-valued 1-form on $P$ with $L_{q}^{*}\omega=\rho(q)\circ \omega$ for all $q\in Q$.
\item
An $\mathfrak{s}$-valued 1-form $\omega$ on $P$ with $\omega(\wt{s}(p))=\vec{s}$ for all $\vec{s}\in
\mathfrak{s}$ is a~$\Phi$-invariant connection if\/f it is of type $\rho$.
\item
Let $Q_{p}$ denote the stabilizer of $p\in P$ w.r.t.~$\Theta$, and $G_{\pi(p)}$ the stabilizer of
$\pi(p)$ w.r.t.~$\varphi$.
Then, $G_{\pi(p)}=\left\{h\in G\: | \:L_h\colon F_{\pi(p)}\rightarrow F_{\pi(p)} \right\}$, and we obtain a~Lie group homomorphism
\begin{gather*}
 \phi_p\colon G_{\pi(p)}\rightarrow S\quad\text{by requiring that}\quad \Phi(h,p)=p\cdot\phi_p(h)\quad\text{for all}\quad h\in
G_{\pi(p)}.
\end{gather*}
If~$\mathfrak{q}_p$ and~$\mathfrak{g}_{\pi(p)}$ denote the Lie algebras of $Q_p$ and $G_{\pi(p)}$, respectively, then
\begin{gather}
\label{eq:staoQ}
Q_p=\{(h,\phi_p(h))\:|\:h\in G_{\pi(p)}\}
\qquad
\text{and}
\qquad
\mathfrak{q}_p=\big\{\big(\vec{h},
\mathrm{d}_e\phi_p\big(\vec{h}\hspace{1pt}\big)\big)\:\big|\:\vec{h}\in
\mathfrak{g}_{\pi(p)}\big\}.
\end{gather}
\end{itemize}


\section[$\Phi$-coverings]{$\boldsymbol{\Phi}$-coverings}\label{sec:phiCoverings}

We start this section with some facts and conventions concerning submanifolds.
Then, we provide the def\/inition of a~$\Phi$-covering and discuss some its properties.
\begin{Convention}
\label{conv:Submnfds}
Let~$M$ be a~manifold.
\begin{enumerate}\itemsep=0pt
\item[$1.$]
A pair $(N,\tau_N)$ consisting of a~manifold $N$ and an injective immersion $\tau\colon N\rightarrow M$
is called submanifold of~$M$.
\item[$2.$]
If $(N,\tau_N)$ is a~submanifold of~$M$, we tacitly identify $N$ and $TN$ with their images
\mbox{$\tau_N(N) \subseteq  M$} and $\mathrm{d}\tau_N(TN)\subseteq TM$, respectively.
In particular, this means that:{\samepage
\begin{itemize}\itemsep=0pt
\item
If $M'$ is a~manifold and $\kappa\colon M\rightarrow M'$ a~smooth map, then for $x\in N$ and
$\vec{v}\in TN$ we write $\kappa(x)$ and $\mathrm{d}\kappa(\vec{v})$ instead of $\kappa(\tau_N(x))$ and
$\mathrm{d}\kappa( \mathrm{d}\tau(\vec{v}))$, respectively.
\item
If $\Psi\colon G\times M\rightarrow M$ is a~left action of the Lie group~$G$ and $(H,\tau_H)$ a~submanifold of~$G$, the restriction of $\Psi$ to $H\times N$ is def\/ined by
\begin{gather*}
\Psi|_{H\times N}(h,x):= \Psi(\tau_H(h),\tau_N(x))
\qquad\forall\, (h,x)\in H\times N.
\end{gather*}
\item
If $\omega\colon TM\rightarrow V$ is a~$V$-valued 1-form on~$M$, we let
\begin{gather*}
(\Psi^{*}\omega)|_{TG\times TN}(\vec{m},\vec{v}):=(\Psi^{*}\omega)(\vec{m},\mathrm{d}\tau(\vec{v}))
\qquad\forall\, (\vec{m},\vec{v})\in TG\times TN.
\end{gather*}
\item
We will not explicitly refer to the maps $\tau_N$ and $\tau_H$ in the following.
\end{itemize}}
\item[$3.$]
Open subsets $U\subseteq M$ are equipped with the canonical manifold structure making the inclusion map an
embedding.
\item[$4.$]
If $L$ is a~submanifold of $N$, and $N$ is a~submanifold of~$M$, we consider $L$ as a~submanifold of~$M$
in the canonical way.
\end{enumerate}
\end{Convention}

\begin{Definition}
A submanifold $N\subseteq M$ is called $\Psi$-patch if\/f for each $x\in N$ we find an open neighbourhood $N'\subseteq N$
of $x$ and a~submanifold $H$ of~$G$ through $e$, such that the restriction $\Psi|_{H\times N'}$ is a~dif\/feomorphism to an
open subset $U\subseteq M$.
\end{Definition}

\begin{Remark}\label{rem:patch}\quad
\begin{enumerate}\itemsep=0pt
\item
It follows from the inverse function theorem and\footnote{The sum is not necessarily direct.}
\begin{gather*}
\mathrm{d}_{(e,x)}\Psi(\mathfrak{g}\times T_xN)= \mathrm{d}_e\Psi_x(\mathfrak{g}) + \mathrm{d}_x\Psi_e(T_xN)=
\mathrm{d}_e\Psi_x(\mathfrak{g})+T_xN
\qquad\forall\, x\in N
\end{gather*}
that $N$ is a~$\Psi$-patch if\/f $T_xM= \mathrm{d}_e\Psi_x(\mathfrak{g})+T_xN$ holds for all $x\in N$.\footnote{In fact, let $V\subseteq \mathrm{d}_e\Psi_x(\mathfrak{g})$ be an algebraic complement of $T_xN$ in $T_xM$ and $V'\subseteq
\mathfrak{g}$ a~linear subspace with $\dim[V']=\dim[V]$ and $\mathrm{d}_e\Psi_x(V')=V$.
Then, we f\/ind a~submanifold $H$ of~$G$ through $e$ with $T_eH=V'$, so that $\mathrm{d}_{(e,x)}\Psi\colon T_eH\times T_xN
\rightarrow T_xM$ is bijective.}
\item
Open subsets $U\subseteq M$ are always $\Psi$-patches.
They are of maximal dimension, which, for instance, is necessary if there is a~point in $U$ whose stabilizer equals~$G$,
see Lemma~\ref{lemma:mindimslice}.1.
\item
We allow zero-dimensional patches, i.e., $N=\{x\}$ for some $x\in M$.
Necessarily, then we have $\mathrm{d}_e\Psi_x(\mathfrak{g})=T_xM$ as well as $\Psi|_{H\times N}=\Psi_x|_{H}$ for each submanifold $H$
of~$G$.
\end{enumerate}
\end{Remark}

The second part of the following elementary lemma equals Lemma~2.1.1 in~\cite{DuisKolk}.
\begin{Lemma}\label{lemma:mindimslice}
Let $(G,\Psi)$ be a~Lie group that acts on the manifold~$M$, and let $x\in M$.
\begin{enumerate}\itemsep=0pt%\label{item:s}
\item[$1.$]
If $N$ is a~$\Psi$-patch with $x\in N$, then $\dim[N]\geq \dim[M]-\dim[G]+\dim[G_x]$.
\item[$2.$]
Let $V$ and $W$ be algebraic complements of $\mathrm{d}_e\Psi_x(\mathfrak{g})$ in $T_xM$ and of $\mathfrak{g}_x$ in
$\mathfrak{g}$, respectively.
Then there are submanifolds $N$ of~$M$ through $x$ and $H$ of~$G$ through $e$ such that $T_xN=V$, $T_eH=W$.
In particular, $N$ is a~$\Psi$-patch and $\dim[N]=\dim[M]-\dim[G]+\dim[G_x]$.
\end{enumerate}
\end{Lemma}

\begin{proof}
1.~By Remark~\ref{rem:patch}.1 and since $\ker[\mathrm{d}_e\Psi_x]=\mathfrak{g}_x$, we have
\begin{gather}
\label{eq:leq}
\dim[M]\leq \dim[\mathrm{d}_e\Psi_x(\mathfrak{g})]+\dim[T_xN]=\dim[G]-\dim[G_x]+\dim[N].
\end{gather}

2.~Of course, we f\/ind submanifolds $N'$ of~$M$ through $x$ and $H'$ of~$G$ through $e$ such that $T_xN'=V$ and $T_eH'=W$.
So, if $\vec{g}\in \mathfrak{g}$ and $\vec{v}_x\in T_xN'$, then
$0=\mathrm{d}_{(e,x)}\Psi(\vec{g},\vec{v}_x)=\mathrm{d}_e\Psi_x(\vec{g})+\vec{v}_x$ implies
$\mathrm{d}_e\Psi_x(\vec{g})=0$ and $\vec{v}_x=0$.
Hence, $\vec{g}\in \ker[\mathrm{d}_e\Psi_x]=\mathfrak{g}_x$, so that\footnote{Recall that
$\mathrm{d}_{(e,x)}\Psi|_{T_eH'\times T_eN'}\colon \big(\vec{h},\vec{v}_x\big)\mapsto
\mathrm{d}_{(e,x)}\Psi\big(\mathrm{d}_e\tau_H(\vec{h}),\mathrm{d}_x\tau_N(\vec{v}_x)\big)$.}
$\mathrm{d}_{(e,x)}\Psi|_{T_eH'\times T_eN'}$ is injective.
It is immediate from the def\/initions that this map is surjective, so that by the inverse function theorem we f\/ind open
neighbourhoods $N\subseteq N'$ of $x$ and $H\subseteq G$ of $e$ such that $\Psi|_{H\times N}$ is a~dif\/feomorphism to an
open subset $U\subseteq M$.
Then $N$ is a~$\Psi$-patch, and since in~\eqref{eq:leq} equality holds, also the last claim is clear.
\end{proof}

\begin{Definition}\looseness=1%\label{def:pmappe}
Let $(G,\Phi)$ be a~Lie group of automorphisms of the principal f\/ibre bundle~$P$, and recall the actions $\varphi$ and
$\Theta$ def\/ined by~\eqref{eq:INDA} and~\eqref{eq:THETA}, respectively.
A family of $\Theta$-patches $\{P_\alpha\}_{\alpha\in I}$ is said to be a~$\Phi$-covering of $P$ if\/f each
$\varphi$-orbit intersects at least one of the sets~$\pi(P_\alpha)$.
\end{Definition}

\begin{Remark}\label{bem:Psliceeigensch}{\samepage
\quad
\begin{enumerate}\itemsep=0pt
\item
If $O\subseteq P$ is a~$\Theta$-patch, Lemma~\ref{lemma:mindimslice}.1 and~\eqref{eq:staoQ} yield
\begin{gather*}
\dim[O]\geq \dim[P]-\dim[Q]+\dim[Q_p] \stackrel{\eqref{eq:staoQ}}{=}\dim[M]-\dim[G]+\dim[G_{\pi(p)}].
\end{gather*}
\item
It follows from Remark~\ref{rem:patch}.1 and
$\mathrm{d}_e\Theta_p(\mathfrak{q})=\mathrm{d}_e\Phi_p(\mathfrak{g}) + Tv_pP$ that $O$ is a~$\Theta$-patch if\/f
\begin{gather}
\label{eq:transv}
T_pP=T_pO+ \mathrm{d}_e\Phi_p(\mathfrak{g}) + Tv_pP
\qquad\forall\, p\in O.
\end{gather}
As a consequence,
\begingroup
\setlength{\leftmarginii}{19pt}
\begin{itemize}\itemsep=0pt
\item
each $\Phi$-patch is a~$\Theta$-patch,
\item
$P$ is always a~$\Phi$-covering by itself. Moreover, if $P=M\times S$ is trivial, then $M\times \{e\}$ is a~$\Phi$-covering.
\end{itemize}
\item
If $N$ is a~$\varphi$-patch and $s_0\colon N \rightarrow P$ a~smooth section (i.e., $\pi\circ s_0
=\mathrm{id}_N$), then $s_0(N)$ is a~$\Theta$-patch by Lemma~\ref{lemma:suralpha}.2.

Conversely, if $N\subseteq M$ is a submanifold such that $s_0(N)$ is a~$\Theta$-patch for $s_0$ as above, then~$N$ is a $\varphi$-patch. In fact, applying $\mathrm{d}\pi$ to \eqref{eq:transv}, this is immediate from Remark \ref{rem:patch}.1 and the definition of $\varphi$.
\end{enumerate}
\endgroup
}
\end{Remark}

\begin{Lemma}\label{lemma:suralpha}
Let $(G,\Phi)$ be a~Lie group of automorphisms of the principal bundle $(P,\pi,M,S)$.
\begin{enumerate}\itemsep=0pt
\item[$1.$]
If $O\subseteq P$ is a~$\Theta$-patch, then for each $p\in O$ and $q\in Q$ the differential
$\mathrm{d}_{(q,p)}\Theta\colon T_qQ\times T_{p}O\rightarrow T_{q\cdot p}P$ is surjective.
\item[$2.$]
If $N$ is a~$\varphi$-patch and $s_0\colon N \rightarrow P$ a~smooth section, then $s_0(N)$ is a~$\Theta$-patch.
\end{enumerate}
\end{Lemma}

\begin{proof}
1.~Since $O$ is a~$\Theta$-patch, the claim is clear for $q=e$.
If $q$ is arbitrary, then for each $\vec{m}_q\in T_qQ$ we f\/ind some $\vec{q}\in\mathfrak{q}$ such that
$\vec{m}_q=\mathrm{d} L_q\vec{q}$.
Consequently, for $\vec{w}_{p}\in T_{p}P$ we have
\begin{gather*}
\mathrm{d}_{(q,p)}\Theta\left(\vec{m}_q,\vec{w}_{p}\right)=\mathrm{d}_{(q,p)}\Theta(\mathrm{d}
L_q\vec{q},\vec{w}_{p})=\mathrm{d}_{p}L_q\left(\mathrm{d}_{(e,p)}\Theta(\vec{q},\vec{w}_{p})\right).
\end{gather*}
So, since left translation w.r.t.~$\Theta$ is a~dif\/feomorphism, $\mathrm{d}_{p}L_q$ is surjective.

2.~$O:=s_0(N)$ is a~submanifold of $P$ because $s_0$ is an injective immersion. Thus,
by Remark~\ref{bem:Psliceeigensch}.2 it suf\/f\/ices to show that
\begin{gather*}
\dim\big[T_{s_0(x)}O + \mathrm{d}_e\Phi_{s_0(x)}(\mathfrak{g})+ Tv_{s_0(x)}P\big]\geq
\dim[T_{s_0(x)}P]
\qquad\forall\, x\in N.
\end{gather*}
For this, let $x\in N$ and $V'\subseteq \mathfrak{g}$ be a~linear subspace with $V'\oplus \mathfrak{g}_x$ and $T_xM=T_xN \oplus
\mathrm{d}_e\varphi_x(V')$.
Then, we have $T_{s_0(x)}O \oplus \mathrm{d}_e\Phi_{s_0(x)}(V')\oplus Tv_{s_0(x)}P$ because if $\mathrm{d}_xs_0(\vec{v}_x)
+\mathrm{d}_e\Phi_{s_0(x)}(\vec{g}\hspace{1pt}{}')+ \vec{v}_v=0$ for $\vec{v}_x\in T_xN$, $\vec{g}\hspace{1pt}{}'\in V'$ and
$\vec{v}_v\in Tv_{s_0(x)}P$,
\begin{gather*}
0=\mathrm{d}_{s_0(x)}\pi \big(\mathrm{d}_xs_0(\vec{v}_x) +\textstyle\mathrm{d}_e\Phi_{s_0(x)}(\vec{g}\hspace{1pt}{}')+
\vec{v}_v\big)=\vec{v}_x \oplus \mathrm{d}_e\varphi_x(\vec{g}\hspace{1pt}{}')
\end{gather*}
shows $\vec{v}_x \!=\!0$ and $\mathrm{d}_e\phi_{x}(\vec{g}')\!=\!0$, hence $\vec{g}'\!=\!0$ by the choice of $V'$, i.e., $\vec{v}_v\!=\!0$ by assumption.
In particular, $\mathrm{d}_e\phi_{x}(\vec{g}')\!=\!0$ if $\mathrm{d}_e\Phi_{s_0(x)}(\vec{g}')\!=\!0$, hence $\dim[\mathrm{d}_e\Phi_{s_0(x)}(V')]$ $\geq\dim[\mathrm{d}_e\varphi_x(V')]$, from which we obtain
\begin{gather*}
\dim\big[T_{s_0(x)}O + \mathrm{d}_e\Phi_{s_0(x)}(\mathfrak{g})+ Tv_{s_0(x)}P\big]
\geq \dim\big[T_{s_0(x)}O \oplus \mathrm{d}_e\Phi_{s_0(x)}(V')\oplus Tv_{s_0(x)}P\big]
\\
\qquad{}
=\dim[T_xN] + \dim[\mathrm{d}_e\Phi_{s_0(x)}(V')] + \dim[S]
\geq\dim[T_xN] + \dim[\mathrm{d}_e\varphi_x(V')] + \dim[S]\\
\qquad{}
=\dim[P].\tag*{\qed}
\end{gather*}
\renewcommand{\qed}{}
\end{proof}

\section{Characterization of invariant connections}\label{sec:mainth}

In this section, we will use $\Phi$-coverings $\{P_\alpha\}_{\alpha\in I}$ of the bundle $P$ in order to characterize the set of
$\Phi$-invariant connections by families $\{\psi_\alpha\}_{\alpha\in I}$ of smooth maps $\psi_\alpha\colon
\mathfrak{g}\times TP_\alpha \rightarrow \mathfrak{s}$ whose restrictions $\psi_\alpha|_{\mathfrak{g}\times
T_{p_\alpha}P_\alpha }$ are linear and that fulf\/il two additional compatibility conditions.
Here, we will follow the lines of Wang's original approach, which basically means that we generalize the proofs from~\cite{Wang} to the
non-transitive case.
We will proceed in two steps, the f\/irst one being performed in Subsection~\ref{subsec:RedInvConn}.
There, we show that a~$\Phi$-invariant connection gives rise to a~consistent family $\{\psi_\alpha\}_{\alpha\in I}$ of
smooth maps as described above.
We also discuss the situation in~\cite{HarSni} in order to make the two conditions more intuitive.
Then, in Subsection~\ref{subsec:Reconstruc}, we will verify that such families $\{\psi_\alpha\}_{\alpha\in I}$ glue together
to a~$\Phi$-invariant connection on $P$.

\subsection{Reduction of invariant connections}\label{subsec:RedInvConn}

In the following, let $\{P_\alpha\}_{\alpha\in I}$ be a~f\/ixed $\Phi$-covering of $P$ and $\omega$ a~$\Phi$-invariant
connection on $P$.
We def\/ine
\begin{gather*}
\omega_\alpha:=(\Theta^{*}\omega)|_{TQ\times TP_\alpha}\qquad\text{as well as}\qquad\psi_\alpha:=\omega_\alpha|_{\mathfrak{g}\times
TP_\alpha},
\end{gather*}
and for $q'\in Q$ we let $\Co{q'}\colon Q\times P\rightarrow Q\times P$, $(q,p)\mapsto \left(\Co{q'}(q),p\right)$. Finally, we define
\begin{gather*}
\Add{q}(\vec{g}):=\Add{g}(\vec{g})\qquad\forall\, q=(g,s)\in Q,\qquad \forall\, \vec{g}\in \mathfrak{g}.
\end{gather*}
\begin{Lemma}
\label{lemma:omegaalpha}
Let $q\in Q$, $p_\alpha\in P_\alpha$, $p_\beta\in P_\beta$ with\footnote{Recall that, by Convention~\ref{conv:Submnfds}, this actually means $\tau_{P_\beta}(p_\beta)=q\cdot
\tau_{P_\alpha}(p_\alpha)$.} $p_\beta=q\cdot p_\alpha$ and $\vec{w}_{p_\alpha}\in
T_{p_\alpha}P_\alpha$.
Then
\begin{enumerate}\itemsep=0pt
\item[$1)$]
$\omega_\beta(\vec{\eta}\hspace{1pt})=\rho(q)\circ \omega_\alpha(\vec{0}_\mathfrak{q},\vec{w}_{p_\alpha})$ for all $\vec{\eta}\in
TQ\times TP_\beta$ with $\mathrm{d}\Theta(\vec{\eta}\hspace{1pt})=\mathrm{d} L_q\vec{w}_{p_\alpha}$,
\item[$2)$]
$\left(\Co{q}^{*}\omega_\beta\right)\big(\vec{m},\vec{0}_{p_\beta}\big)=\rho(q)\circ
\omega_\alpha\big(\vec{m},\vec{0}_{p_\alpha}\big)$ for all $\vec{m}\in TQ$.
\end{enumerate}
\end{Lemma}

\begin{proof}
1.~Let $\vec{\eta}\in T_{q'}Q\times T_pP_\beta$ for $q'\in Q$.
Then, since\footnote{See end of Subsection~\ref{subsec:InvConn}.} $L_q^{*}\omega = \rho(q)\circ \omega$ for each $q\in Q$
and $q'\cdot p= q\cdot p_\alpha =p_\beta$, we have
\begin{gather*}
\begin{split}
&\omega_\beta(\vec{\eta}\hspace{1pt})=\omega_{q'\cdot p}(\mathrm{d}_{(q', p)}\Theta (\vec{\eta}\hspace{1pt}))=\omega_{p_\beta}(\mathrm{d}
L_q\vec{w}_{p_\alpha}) =(L_q^{*}\omega)_{p_\alpha}(\vec{w}_{p_\alpha})
\\
&\phantom{\omega_\beta(\vec{\eta})}=\rho(q)\circ \omega_{p_\alpha}(\vec{w}_{p_\alpha})=\rho(q)\circ \omega_{p_\alpha}\big(\mathrm{d}_{(e,p_\alpha)}\Theta
\big(\vec{0}_\mathfrak{q},\vec{w}_{p_\alpha}\big)\big)
=\rho(q)\circ
\omega_{\alpha}\big(\vec{0}_\mathfrak{q},\vec{w}_{p_\alpha}\big).
\end{split}
\end{gather*}

2.~For $\vec{m}_{q'} \in T_{q'}Q$ let $\gamma \colon (-\epsilon,\epsilon)\rightarrow Q$ be smooth with
$\dot\gamma(0)=\vec{m}_{q'}$.
Then
\begin{gather*}
\big(\Co{q}^{*}\omega_\beta
\big)_{(q',p_\beta)}\big(\vec{m}_{q'},\vec{0}_{p_\beta}\big)
={\omega_\beta}_{(\Co{q}(q'),p_\beta)}\big(\Add{q}(\vec{m}_{q'}),\vec{0}_{p_\beta}\big)
=\omega_{qq'q^{-1}q\cdot p_\alpha }\left(\dttB{t}{0}q\gamma(t) q^{-1}q\cdot p_\alpha\right)
\\
\hphantom{\big(\Co{q}^{*}\omega_\beta
\big)_{(q',p_\beta)}\big(\vec{m}_{q'},\vec{0}_{p_\beta}\big)}{}
=\left(L_q^{*}\omega\right)_{q'\cdot p_\alpha }\left(\dttB{t}{0}\gamma(t)\cdot p_\alpha\right)
=\rho(q)\circ \omega_{q'\cdot p_\alpha}\left(\mathrm{d}_{(q',p_\alpha)}\Theta\left( \vec{m}_{q'}\right)\right)
\\
\hphantom{\big(\Co{q}^{*}\omega_\beta
\big)_{(q',p_\beta)}\big(\vec{m}_{q'},\vec{0}_{p_\beta}\big)}{}
=\rho(q)\circ {\omega_\alpha}_{(q',p_\alpha)} \big(\vec{m}_{q'},\vec{0}_{p_\alpha}\big).\tag*{\qed}
\end{gather*}
\renewcommand{\qed}{}
\end{proof}


\begin{Corollary}
\label{cor:psialpha}
Let $q\in Q$, $p_\alpha\in P_\alpha$, $p_\beta\in P_\beta$ with $p_\beta=q\cdot p_\alpha$ and $\vec{w}_{p_\alpha}\in
T_{p_\alpha}P_\alpha$.
Then, for $\vec{w}_{p_\beta}\in T_{p_\beta}P_\beta$, $\vec{g}\in \mathfrak{g}$ and $\vec{s}\in \mathfrak{s}$ we have
\begin{enumerate}\itemsep=0pt
\item[$i)$] $\wt{g}(p_\beta) + \vec{w}_{p_\beta}-\wt{s}(p_\beta)=\mathrm{d} L_q\vec{w}_{p_\alpha}
\quad
\Longrightarrow
\quad
\psi_\beta(\vec{g},\vec{w}_{p_\beta})-\vec{s}
=\rho(q)\circ\psi_\alpha\big(\vec{0}_{\mathfrak{g}},\vec{w}_{p_\alpha}\big)$,
\item[$ii)$] $\psi_\beta\big(\Add{q}(\vec{g}),\vec{0}_{p_\beta}\big)=\rho(q)\circ
\psi_\alpha\big(\vec{g},\vec{0}_{p_\alpha}\big)$.
\end{enumerate}
\end{Corollary}

\begin{proof}
$i)$ In general, for $\vec{w}_p\in T_pP$, $\vec{g}\in\mathfrak{g}$ and $\vec{s}\in\mathfrak{s}$ we have
\begin{gather}
\label{eq:ThetaPhi}
\mathrm{d}_{(e,p)}\Theta((\vec{g},\vec{s}\hspace{1pt}),\vec{w}_p)
=\mathrm{d}_{(e,p)}\Phi(\vec{g},\vec{w}_p)-\wt{s}(p)=\wt{g}(p)+\vec{w}_p-\wt{s}(p)
\end{gather}
and, since $\omega$ is a~connection, for $((\vec{g},\vec{s}),\vec{w}_{p_\alpha})\in \mathfrak{q}\times TP_\alpha$ we
obtain{\samepage
\begin{gather}
\begin{split}
& \omega_\alpha((\vec{g},\vec{s}\hspace{1pt}),\vec{w}_{p_\alpha})
=\omega\big(\mathrm{d}_{(e,p_\alpha)}\Phi(\vec{g},\vec{w}_{p_\alpha})-\wt{s}(p_\alpha)\big)
=\omega\big(\mathrm{d}_{(e,p_\alpha)}\Phi(\vec{g},\vec{w}_{p_\alpha})\big)-\vec{s}
\\
&\phantom{\omega_\alpha((\vec{g},\vec{s}),\vec{w}_{p_\alpha})}
=\omega_\alpha\left(\vec{g},\vec{w}_{p_\alpha}\right)-\vec{s}
=\psi_\alpha\left(\vec{g},\vec{w}_{p_\alpha}\right)-\vec{s}.
\end{split} \label{eq:gugu}
\end{gather}
Now,} assume that $\mathrm{d}_e\Phi_{p_\beta}(\vec{g}\hspace{1pt})+\vec{w}_{p_\beta}-\wt{s}(p)=\mathrm{d}
L_q\vec{w}_{p_\alpha}$.
Then $\mathrm{d}_{(e,p_\beta)}\Theta((\vec{g},\vec{s}\hspace{1pt}),\vec{w}_{p_\beta})=\mathrm{d} L_q \vec{w}_{p_\alpha}$
by~\eqref{eq:ThetaPhi} so that $\omega_\beta((\vec{g},\vec{s}),\vec{w}_{p_\beta})=\rho(q)\circ
\omega_\alpha\big(\vec{0}_{\mathfrak{g}},\vec{w}_{p_\alpha}\big)$ by Lemma~\ref{lemma:omegaalpha}.1.
Consequently,
\begin{gather*}
\psi_\beta\left(\vec{g},\vec{w}_{p_\beta}\right)-\vec{s}
\stackrel{\eqref{eq:gugu}}{=}\omega_\beta((\vec{g},\vec{s}\hspace{1pt}),\vec{w}_{p_\beta})
=\rho(q)\circ
\omega_\alpha\big(\vec{0}_\mathfrak{q},\vec{w}_{p_\alpha}\big)\stackrel{\eqref{eq:gugu}}{=}\rho(q)\circ
\psi_\alpha\big(\vec{0}_{\mathfrak{g}},\vec{w}_{p_\alpha}\big).
\end{gather*}

$ii)$ Lemma~\ref{lemma:omegaalpha}.2 yields
\begin{gather*}
\psi_\beta\big(\Add{q}(\vec{g}\hspace{1pt}),\vec{0}_{p_\beta}\big)
=(\Co{q}^{*}\omega_\beta)_{(e,p_\beta)}\big(\vec{g},\vec{0}_{p_\beta}\big)
=\rho(q)\circ (\omega_\alpha)_{(e,p_\alpha)}\big(\vec{g},\vec{0}_{p_\alpha}\big) =\rho(q)\circ
\psi_\alpha\big(\vec{g},\vec{0}_{p_\alpha}\big).\!\!\!\!\tag*{\qed}
\end{gather*}
\renewcommand{\qed}{}
\end{proof}


\begin{Definition}[reduced connection]
A family $\{\psi_\alpha\}_{\alpha\in I}$ of smooth maps $\psi_\alpha\colon \mathfrak{g}\times TP_\alpha \rightarrow
\mathfrak{s}$ which are linear in the sense that $\psi_\alpha|_{\mathfrak{g}\times T_{p_\alpha}P_\alpha }$ is linear for
all $p_\alpha\in P_\alpha$ is called reduced connection w.r.t.~$\{P_\alpha\}_{\alpha\in I}$ if\/f it fulf\/ils the
conditions $i)$ and $ii)$ from Corollary~\ref{cor:psialpha}.
\end{Definition}
\begin{Remark}\label{rem:Psialphaconderkl}\quad
\begin{enumerate}\itemsep=0pt
\item[1)] In particular, Corollary~\ref{cor:psialpha}.$i)$ encodes the following condition
\begin{enumerate}\itemsep=0pt
\item[$a)$] For all $\beta \in I$, $(\vec{g},\vec{s}\hspace{1pt})\in \mathfrak{q}$ and $\vec{w}_{p_\beta}\in
T_{p_\beta}P_\beta$ we have
\begin{gather*}
\wt{g}(p_\beta) +\vec{w}_{p_\beta}-\wt{s}(p_\beta)=0
\quad
\Longrightarrow
\quad
\psi_\beta(\vec{g},\vec{w}_{p_\beta})-\vec{s}=0.
\end{gather*}
\end{enumerate}
\item[2)] Assume that $a)$ is true and let $q\in Q$, $p_\alpha\in P_\alpha$, $p_\beta \in
P_\beta$ with $p_\beta=q\cdot p_\alpha$.
Moreover, assume that we f\/ind elements $\vec{w}_{p_\alpha}\in T_{p_\alpha}P_\alpha$ and
$((\vec{g},\vec{s}\hspace{1pt}),\vec{w}_{p_\beta})\in \mathfrak{q}\times T_{p_\beta}P_\beta$ such that
\begin{gather*}
\mathrm{d}_{(e,p_\beta)}\Theta((\vec{g},\vec{s}\hspace{1pt}),\vec{w}_{p_\beta})=\mathrm{d} L_q\vec{w}_{p_\alpha}
\qquad
\text{and}
\qquad
\psi_\beta(\vec{g},\vec{w}_{p_\beta})-\vec{s}=\rho(q)\circ\psi_\alpha(\vec{0}_{\mathfrak{g}},\vec{w}_{p_\alpha})
\end{gather*}
holds.
Then
$\psi_\beta\big(\vec{g}\hspace{1pt}{}',\vec{w}'_{p_\beta}\big)-\vec{s}\hspace{1pt}{}'
=\rho(q)\circ\psi_\alpha\big(\vec{0}_{\mathfrak{g}},\vec{w}_{p_\alpha}\big)$
holds for each element\footnote{Observe that due to surjectivity of $\mathrm{d}_{(e,p_\beta)}\Phi$ such elements always exist.}
$\big((\vec{g}\hspace{1pt}{}',\vec{s}\hspace{1pt}{}'),\vec{w}'_{p_\beta}\big)\in \mathfrak{q}\times T_{p_\beta}P_\beta$
with\footnote{Recall equation~\eqref{eq:ThetaPhi}.}
$\mathrm{d}_{(e,p_\beta)}\Theta\big((\vec{g}\hspace{1pt}{}',\vec{s}\hspace{1pt}{}'),\vec{w}'_{p_\beta}\big)=\mathrm{d}
L_q\vec{w}_{p_\alpha}$.
In fact, we have
\begin{gather*}
\mathrm{d}_{(e,p_\beta)}\Theta\big((\vec{g}-\vec{g}\hspace{1pt}{}',\vec{s}-\vec{s}\hspace{1pt}{}'),\vec{w}_{p_\beta}-\vec{w}'_{p_\beta}\big)=0,
\end{gather*}
so that by~\eqref{eq:ThetaPhi} condition $a)$ gives
\begin{gather*}
0\stackrel{a)}{=}\psi_\beta(\vec{g}-\vec{g}\hspace{1pt}{}',\vec{w}_{p_\beta}-\vec{w}'_{p_\beta})-(\vec{s}-\vec{s}\hspace{1pt}{}'))
=\big[\psi_\beta(\vec{g},\vec{w}_{p_\beta})-\vec{s}\big]-\big[\psi_\beta(\vec{g}\hspace{1pt}{}',\vec{w}'_{p_\beta})-\vec{s}\hspace{1pt}{}'\big]
\\
\phantom{0}=
\rho(q)\circ\psi_\alpha\big(\vec{0}_{\mathfrak{g}},\vec{w}_{p_\alpha}\big)
-\big[\psi_\beta(\vec{g}\hspace{1pt}{}',\vec{w}'_{p_\beta})-\vec{s}\hspace{1pt}{}'\big].
\end{gather*}
\item[3)] Assume that $\mathrm{d} L_q \vec{w}_{p_\alpha}\in T_{p_\beta}P_\beta$ holds for all $q\in Q$,
$p_\alpha\in P_\alpha$, $p_\beta \in P_\beta$ with $p_\beta=q\cdot p_\alpha$ and all $\vec{w}_{p_\alpha}\in
T_{p_\alpha}P_\alpha$.
Then $\mathrm{d}_{(e,p_\beta)}\Theta\left({\rm d}L_q \vec{w}_{p_\alpha}\right)=\mathrm{d} L_q \vec{w}_{p_\alpha}$ so that it
follows from 2) that in this case we can substitute $i)$ by $a)$ and condition
\begin{enumerate}\itemsep=0pt
\item[$b)$] Let $q\in Q$, $p_\alpha\in P_\alpha$, $p_\beta \in P_\beta$ with $p_\beta=q\cdot
p_\alpha$.
Then
\begin{gather*}
\psi_\beta\big(\vec{0}_{\mathfrak{g}},\mathrm{d}
L_q\vec{w}_{p_\alpha}\big)=\rho(q)\circ\psi_\alpha\big(\vec{0}_{\mathfrak{g}},\vec{w}_{p_\alpha}\big)
\qquad\forall\, \vec{w}_{p_\alpha}\in T_{p_\alpha}P_\alpha.
\end{gather*}
\end{enumerate}
Now, $b)$ looks similar to $ii)$ and makes it plausible that the conditions
$i)$ and $ii)$ from Corollary~\ref{cor:psialpha} together encode the $\rho$-invariance of the corresponding
connection $\omega$.
However, usually there is no reason for $\mathrm{d} L_q \vec{w}_{p_\alpha}$ to be an element of $T_{p_\beta}P_\beta$.
Even for $p_\alpha=p_\beta$ and $q\in Q_{p_\alpha}$ this is usually not true.
Thus, typically there is no way to split up $i)$ into parts whose meaning is more intuitive.
\end{enumerate}
\end{Remark}
